\chapter{Einleitung}
\label{ch:einleitung}

Viele geschäftskritische Legacy-Systeme werden weit über ihren geplanten Lebenszyklus hinaus betrieben und weiterentwickelt. Obwohl ihre Programmstrukturen und Technologien oft veraltet sind, bleiben sie in produktiven Umgebungen unverzichtbar~\citep{363157}. Ihre Weiterentwicklung konzentriert sich meist auf funktionale Anforderungen, während Wartbarkeit vernachlässigt werden~\citep{refactoring_for_software_architecture_smells}. Dadurch entstehen sogenannte Design Smells. Design Smells sind keine direkten Programmfehler, sondern Hinweise auf strukturelle Probleme, die die Verständlichkeit, Wartbarkeit und Testbarkeit eines Systems verschlechtern~\citep{software_design_smell_detection}. Sie entstehen oft durch bewusste Abkürzungen zugunsten einer schnelleren Auslieferung und erschweren die weitere Entwicklung des Systems~\citep{sharma2020}. Um solche Probleme zu beheben, werden \textit{Refactoring} und \textit{Reengineering} eingesetzt~\citep{laguna_systematic_2013, majthoub_software_2018}. Während \textit{Refactoring} die interne Codestruktur verbessert, kann \textit{Reengineering} auch die Architektur eines Systems verändern~\citep{laguna_systematic_2013}. Diese Masterarbeit betrachtet verteilte Legacy-Systeme, die langlebig, geschäftskritisch und im Produktivbetrieb sind. Ihre Codebasis ist umfangreich, aber kaum dokumentiert, und die ursprünglichen Entwickler stehen nicht mehr zur Ver\-fü\-gung. Dadurch wird das interne Systemverhalten nicht mehr im Detail verstanden und jede Code-Änderung kann unvorhersehbare Seiteneffekte auslösen. Aus diesem Grund konzentriert sich diese Masterarbeit auf ein \textit{Reengineering} der Architektur, das die Auswirkungen der Design Smells reduziert, ohne den Quellcode des Legacy-Systems zu ändern. Ziel ist es, die Architektur an aktuelle Anforderungen an Wartbarkeit und Antwortzeiten anzupassen, während die bestehende Funktionalität erhalten bleibt.

\section{Aufbau der Arbeit}
\label{sec:aufbau}

Zunächst werden in Kapitel~\ref{ch:grundlagen} die Grundlagen dargestellt, die zum Verständnis dieser Masterarbeit relevant sind. Danach wird in Kapitel~\ref{ch:fallstudie} ein verteiltes Legacy-System als Fallstudie betrachtet. Dabei werden die Architektur des Systems beschrieben und dessen Design Smells identifiziert. Anschließend wird das \acrfull{lps}~\citep{liebig_2025_architectural_refactoring} als ein \textit{Architektur-Reengineering-Ansatz} vorgestellt, das die Auswirkungen dieser Design Smells mitigiert. Durch gezielte Experimente wird gezeigt, dass \gls{lps} unter bestimmten Fehlerbedingungen die Delivery-Garantien \textit{Exactly-Once-Delivery} und \textit{In-Order-Delivery} verletzt. Diese Limitationen motivieren die Entwicklung von \acrfull{alps}, das in Kapitel~\ref{ch:alps} vorgestellt wird. \gls{alps} behält die Vorteile von \gls{lps} bei, adressiert jedoch dessen Limitationen. Dafür werden zunächst die Anforderungen an \gls{alps} abgeleitet und anschließend dessen Architektur, Integration in das Legacy-System und Implementierung beschrieben. Abschließend wird \gls{alps} experimentell evaluiert, um die Einhaltung der Delivery-Garantien sowohl unter Fehlerbedingungen als auch unter realitätsnahen Bedingungen zu überprüfen. Kapitel~\ref{ch:fazit} fasst die Ergebnisse dieser Masterarbeit zusammen und Kapitel~\ref{ch:ausblick} gibt einen Ausblick auf mögliche weiterführende Arbeiten.